\subsection{二面体群}

定义 2. 二面体群是由两个不同的二阶元素生成的群。

例。设 M 为乘法群 \(\{1,-1\}\)，并设 m 是一个 \(\geqslant2\) 或 \(m=\infty\) 的整数。于是 M 通过 \((-1).x=-x\) 作用于群 Z/mZ（分别地，作用于 Z），相应的 M 与 Z/mZ（分别地，M 与 Z）的半直积记为 \(D_{m}\)。因此，\(D_{m}\) 的元素是形如 \((\varepsilon,x)\) 的对，其中 \(\varepsilon=\pm1\) 且 \(x\in Z/mZ\)（分别地，\(x\in Z\)）；\(D_{m}\) 上的群运算由公式

\[
(\varepsilon , x). (\varepsilon^ {\prime}, x ^ {\prime}) = (\varepsilon \varepsilon^ {\prime}, \varepsilon^ {\prime} x + x ^ {\prime}).\tag{6}
\]

我们把 \(\iota\) 记作模 \(m\) 的 1 的剩余类（分别地，\(\iota = 1\)），并置

\[
\rho = (- 1, 0), \quad \rho^ {\prime} = (- 1, \iota), \quad \pi = (1, \iota).\tag{7}
\]

于是 \(\rho^2 = \rho'^2 = 1\) 且 \(\pi = \rho \rho'\)。公式

\[
\pi^ {n} = (1, n \iota), \quad \rho \pi^ {n} = (- 1, n \iota)\tag{8}
\]

表明 \(D_{m}\) 是由 \(\{\rho,\rho'\}\) 生成的二面体群。

命题 2. 假设 \(S\) 由两个不同的二阶元素 \(s\) 和 \(s'\) 组成。

\begin{enumerate}
\def\labelenumi{(\roman{enumi})}
\item
  由 \(p = ss'\) 生成的 W 的子群 P 是正规子群，且 W 是子群 \(T = \{1, s\}\) 与 P 的半直积。此外，\((W : P) = 2\)。
\item
  令 \(m\) 为 \(p\) 的阶（有限或无限）。则 \(m \geqslant 2\)，且 \(W\) 的阶为 \(2m\)。存在唯一同构 \(\varphi\)，从 \(D_m\) 到 \(W\)，满足 \(\varphi(\rho) = s\) 且 \(\varphi(\rho') = s'\)。
\item
  我们有 \(sps^{-1} = sss's = s's = p^{-1}\)，因而
\end{enumerate}

\[
s p ^ {n} s ^ {- 1} = p ^ {- n}\tag{9}
\]

对每个整数 \(n\) 都成立。由于 \(W\) 由 \(\{s, s'\}\) 生成，因而也由 \(\{s, p\}\) 生成，子群 \(P\) 是正规子群。因此 \(TP\) 是 \(W\) 的一个子群；由于 \(TP\) 包含 \(s\) 和 \(s' = sp\)，有 \(W = TP = P \cup sP\)。为了证明 (i)，因此只需证明 \(W \neq P\)。若 \(W = P\)，则群 \(W\) 是交换群，所以 \(p^2 = s^2 s'^2 = 1\)。但此时 \(W = P\) 中只有 1 和 \(p\) 两个元素，这与 \(W\) 至少包含三个元素（即 1、\(s\) 和 \(s'\)）的假设矛盾。

\begin{enumerate}
\def\labelenumi{(\roman{enumi})}
\setcounter{enumi}{1}
\tightlist
\item
  由于 \(s \neq s'\)，有 \(p \neq 1\)，故 \(m \geqslant 2\)。由于 P 的阶为 m 且 (W : P) = 2，W 的阶为 2m。当 m 有限（分别地，无限）时，存在一个同构 \(\varphi'\)，从 Z/mZ（分别地，Z）到 P，将 \(\pi\) 映到 p。~此外，存在一个同构 \(\varphi''\)，从 M = \{1, -1\} 到 T，将 -1 映到 s。群 W 是 T 与 P 的半直积。根据公式 (9) 以及 \(\rho\pi^{n}\rho^{-1} = \pi^{-n}\)，\(\varphi'\) 和 \(\varphi''\) 诱导出同构 \(\varphi\)，从 D\_\{m\} 到 W，满足 \(\varphi(\rho) = s\) 和 \(\varphi(\pi) = p\)，从而 \(\varphi(\rho') = s'\)。\(\varphi\) 的唯一性由 D\_\{m\} 由 \(\{\rho, \rho'\}\) 生成这一事实得到。
\end{enumerate}

备注。考虑一个阶为 2m、由两个不同的二阶元素 s 和 \(s'\) 生成的二面体群 W。记 \(s_{q}\)（分别地，\(s'_{q}\)）为长度为 q 的序列，其编号为奇数（分别地，偶数）的项等于 s，而编号为偶数（分别地，奇数）的项等于 \(s'\)；并记 \(w_{q}\)（分别地，\(w'_{q}\)）为序列 \(s_{q}\)（分别地，\(s'_{q}\)）的乘积。我们有

\[
\begin{array}{r l} w _ {2 k} = (s s ^ {\prime}) ^ {k}, & w _ {2 k + 1} = (s s ^ {\prime}) ^ {k} s, \\ w _ {2 k} ^ {\prime} = (s ^ {\prime} s) ^ {k} = (s s ^ {\prime}) ^ {- k}, & w _ {2 k + 1} ^ {\prime} = (s ^ {\prime} s) ^ {k} s ^ {\prime} = (s s ^ {\prime}) ^ {- k - 1} s. \end{array}
\]

如果 \(\mathbf{s} = (s_1, \ldots, s_q)\) 是关于 \(\{s, s'\}\) 的 W 中元素 \(w\) 的一个约化分解，那么显然 \(s_i \neq s_{i+1}\)，其中 \(1 \leqslant i \leqslant q-1\)。因此，\(\mathbf{s} = \mathbf{s}_q\) 或 \(\mathbf{s} = \mathbf{s}_q'\)。

若 \(m = \infty\)，则 \((ss')^n\) 和 \((ss')^ns\) 对于 \(n \in \mathbf{Z}\) 彼此不同。因此，\(w_q\)（\(q \geqslant 0\)）和 \(w_q'\)（\(q > 0\)）彼此不同；若 \(\mathbf{s}\) 是 \(w_q\)（分别地，\(w_q'\)）的一个约化分解，则必有 \(\mathbf{s} = \mathbf{s}_q\)（分别地，\(\mathbf{s} = \mathbf{s}_q'\)）。由此可得 \(l(w_q) = l(w_q') = q\)，且 \(\mathbf{W}\) 中元素的约化分解集合由 \(\mathbf{s}_q\) 和 \(\mathbf{s}_q'\) 组成。此外，\(\mathbf{W}\) 的每个元素都有唯一的约化分解。

现在设 \(m\) 有限。若 \(q \geqslant 2m\)，则 \(w_{q} = w_{q - 2m}\) 且 \(w_{q}' = w_{q - 2m}'\)；若 \(m \leqslant q \leqslant 2m\)，则 \(w_{q} = w_{2m - q}'\)，\(w_{q}' = w_{2m - q}\)。因此，\(\mathbf{s}_q\) 和 \(\mathbf{s}_q'\) 在 \(q > m\) 时都不是约化分解。由此可知，\(2m\) 个元素 \(W\) 中的每一个，都是以下 \(2m\) 个元素之一：\(w_0 = w_0'\)、\(w_q\) 和 \(w_q'\)（其中 \(1 \leqslant q \leqslant m - 1\)），以及 \(w_m = w_m'\)。因此这 \(2m\) 个元素彼此不同，并且如上所述，\(l(w_q) = l(w_q') = q\) 当 \(q \leqslant m\) 时成立，而 \(W\) 中元素的约化分解集合由 \(s_q\) 和 \(s_q'\)（当 \(0 \leqslant q \leqslant m\) 时）组成。\(W\) 中除 \(w_m\) 外的每个元素都有唯一的约化分解；\(w_m\) 有两个。

